\documentclass[11pt]{amsart}
\usepackage{amsmath,amssymb,amsthm}
\usepackage{hyperref}

\newtheorem{theorem}{Theorem}[section]
\newtheorem{proposition}[theorem]{Proposition}
\newtheorem{lemma}[theorem]{Lemma}
\newtheorem{fact}[theorem]{Fact}
\theoremstyle{definition}
\newtheorem{definition}[theorem]{Definition}
\theoremstyle{remark}
\newtheorem{remark}[theorem]{Remark}

\newcommand{\Z}{\mathbb{Z}}
\newcommand{\Q}{\mathbb{Q}}
\newcommand{\R}{\mathbb{R}}
\newcommand{\C}{\mathbb{C}}
\newcommand{\OK}{\mathcal{O}_K}
\newcommand{\Cl}{\operatorname{Cl}}
\newcommand{\Nm}{\operatorname{N}}
\newcommand{\vol}{\operatorname{vol}}
\newcommand{\covol}{\operatorname{covol}}

\title{Finiteness of the class group of a number field}
\date{\today}

\begin{document}
\maketitle

\section{Setting and notation}\label{sec:setting}

Throughout, $K$ is a number field of degree $n = [K:\Q]$, with ring of
integers $\OK$ and discriminant $d_K$. Let $\sigma_1, \dots, \sigma_{r_1}$ be
the real embeddings of $K$ and
$\sigma_{r_1+1}, \overline{\sigma}_{r_1+1}, \dots, \sigma_{r_1+r_2},
\overline{\sigma}_{r_1+r_2}$ the complex ones, so that $n = r_1 + 2r_2$. The
\emph{Minkowski embedding} is
\[
  j \colon K \to \R^{r_1} \times \C^{r_2} \cong \R^n, \qquad
  j(\alpha) = \bigl(\sigma_1(\alpha), \dots, \sigma_{r_1+r_2}(\alpha)\bigr),
\]
where $\C \cong \R^2$ via $z \mapsto (\operatorname{Re} z, \operatorname{Im} z)$
and $\R^n$ carries the Lebesgue measure. For $\alpha \in K$ we write
$\Nm(\alpha) = \Nm_{K/\Q}(\alpha) = \prod_{i=1}^{n} \tau_i(\alpha)$, the
product over all $n$ embeddings $\tau_i$ of $K$ into $\C$.

\begin{definition}\label{def:norm}
For a nonzero ideal $I \subseteq \OK$, its \emph{norm} is
$\Nm(I) = |\OK / I|$.
\end{definition}

\begin{definition}\label{def:classgroup}
A \emph{fractional ideal} of $K$ is a nonzero $\OK$-submodule $\mathfrak{a}$ of
$K$ such that $\delta\mathfrak{a} \subseteq \OK$ for some nonzero
$\delta \in \OK$. The \emph{class group} $\Cl_K$ is the group of fractional
ideals modulo the subgroup of principal fractional ideals $\alpha\OK$,
$\alpha \in K^\times$. For a fractional ideal $\mathfrak{a}$ we write
$[\mathfrak{a}]$ for its class.
\end{definition}

The aim is the following theorem.

\begin{theorem}\label{thm:finite}
The class group $\Cl_K$ is finite.
\end{theorem}

\section{Facts used}\label{sec:facts}

\begin{fact}\label{fact:dedekind}
The ring $\OK$ is a Dedekind domain. Consequently the fractional ideals of $K$
form an abelian group under multiplication, with identity $\OK$, so that
$\Cl_K$ in Definition~\ref{def:classgroup} is a group; and for nonzero ideals
$I, J \subseteq \OK$ we have $J \subseteq I$ if and only if $I \mid J$, that
is, $J = I I'$ for some ideal $I' \subseteq \OK$ \cite[Ch.~3]{Mar}.
\end{fact}

\begin{fact}\label{fact:norm}
The norm of ideals is multiplicative, $\Nm(IJ) = \Nm(I)\Nm(J)$, and for
nonzero $\alpha \in \OK$ we have $\Nm(\alpha\OK) = |\Nm(\alpha)|$
\cite[Ch.~3]{Mar}, \cite[Ch.~I, \S6]{Neu}.
\end{fact}

\begin{fact}\label{fact:covolume}
For every nonzero ideal $I \subseteq \OK$, $j(I)$ is a lattice of rank $n$ in
$\R^n$, and
\[
  \covol\bigl(j(\OK)\bigr) = 2^{-r_2} \sqrt{|d_K|}, \qquad
  \covol\bigl(j(I)\bigr) = \Nm(I) \covol\bigl(j(\OK)\bigr).
\]
See \cite[Ch.~I, \S5]{Neu}; there the covolumes are computed for a measure
equal to $2^{r_2}$ times the Lebesgue measure, which accounts for the factor
$2^{-r_2}$.
\end{fact}

\begin{fact}[Minkowski's lattice point theorem]\label{fact:minkowski}
Let $\Lambda \subseteq \R^n$ be a lattice and $X \subseteq \R^n$ a convex,
centrally symmetric ($x \in X \Rightarrow -x \in X$) set with
$\vol(X) > 2^n \covol(\Lambda)$. Then $X$ contains a nonzero point of
$\Lambda$ \cite[Ch.~I, \S4]{Neu}.
\end{fact}

\begin{lemma}\label{lem:minkowski-compact}
Let $\Lambda \subseteq \R^n$ be a lattice and $X \subseteq \R^n$ a compact,
convex, centrally symmetric set with $\vol(X) \geq 2^n \covol(\Lambda)$. Then
$X$ contains a nonzero point of $\Lambda$.
\end{lemma}

\begin{proof}
Since $X$ is nonempty (its volume is positive) and centrally symmetric, it
contains some $x$ and $-x$, hence $0 = \tfrac12 x + \tfrac12(-x)$ by
convexity. Consequently $sX \subseteq tX$ for $0 \leq s \leq t$, because
$sx = \tfrac{s}{t}(tx) + (1 - \tfrac{s}{t}) \cdot 0$.

For $k \geq 1$ let $X_k = (1 + \tfrac1k)X$. It is convex and centrally
symmetric, and $\vol(X_k) = (1 + \tfrac1k)^n \vol(X) > 2^n \covol(\Lambda)$,
since $\vol(X) \geq 2^n\covol(\Lambda) > 0$. By Fact~\ref{fact:minkowski}
there is a nonzero $\lambda_k \in \Lambda \cap X_k$, and
$X_k \subseteq 2X$ by the previous paragraph. The set $2X$ is bounded and
$\Lambda$ is discrete, so $\Lambda \cap 2X$ is finite. Hence some nonzero
$\lambda \in \Lambda$ equals $\lambda_k$ for infinitely many $k$. For each
such $k$, $\lambda/(1 + \tfrac1k) \in X$; letting $k \to \infty$ along these
indices and using that $X$ is closed, we get $\lambda \in X$.
\end{proof}

\begin{lemma}\label{lem:norm-in-ideal}
For every nonzero ideal $I \subseteq \OK$, the norm $\Nm(I)$ is finite and
$\Nm(I) \in I$.
\end{lemma}

\begin{proof}
By Fact~\ref{fact:covolume}, $j(I)$ and $j(\OK)$ are lattices of rank $n$
with $j(I) \subseteq j(\OK)$, so the index $|\OK / I|$ is finite. Let
$m = \Nm(I) = |\OK/I|$. By Lagrange's theorem applied to the additive group
$\OK/I$ of order $m$, we have $m \cdot (1 + I) = 0$ in $\OK/I$, that is,
$m \in I$.
\end{proof}

\section{The Minkowski bound}\label{sec:bound}

Set
\[
  M_K = \Bigl(\frac{2}{\pi}\Bigr)^{r_2} \sqrt{|d_K|}.
\]

\begin{lemma}\label{lem:minkowski-bound}
Every nonzero ideal $I \subseteq \OK$ contains a nonzero $\alpha$ with
$|\Nm(\alpha)| \leq M_K \Nm(I)$.
\end{lemma}

\begin{proof}
Fix real numbers $c_1, \dots, c_{r_1+r_2} > 0$ with
\begin{equation}\label{eq:product}
  \prod_{i=1}^{r_1} c_i \prod_{i=r_1+1}^{r_1+r_2} c_i^2 = M_K \Nm(I),
\end{equation}
and let
\[
  X_c = \bigl\{ (y_1, \dots, y_{r_1+r_2}) \in \R^{r_1} \times \C^{r_2} :
  |y_i| \leq c_i \text{ for all } i \bigr\}.
\]
This is a product of closed intervals $[-c_i, c_i]$ and closed discs of radius
$c_i$; hence it is compact, convex and centrally symmetric, and by Fubini's
theorem
\[
  \vol(X_c) = \prod_{i=1}^{r_1} 2c_i \prod_{i=r_1+1}^{r_1+r_2} \pi c_i^2
  = 2^{r_1} \pi^{r_2} M_K \Nm(I)
  = 2^{r_1 + r_2} \sqrt{|d_K|}\, \Nm(I),
\]
using \eqref{eq:product} and the definition of $M_K$. On the other hand, by
Fact~\ref{fact:covolume} and $n = r_1 + 2r_2$,
\[
  2^n \covol\bigl(j(I)\bigr) = 2^{n - r_2} \sqrt{|d_K|}\, \Nm(I)
  = 2^{r_1 + r_2} \sqrt{|d_K|}\, \Nm(I).
\]
So $\vol(X_c) = 2^n \covol(j(I))$, and Lemma~\ref{lem:minkowski-compact}
applied to
$\Lambda = j(I)$ and $X = X_c$ gives a nonzero $\alpha \in I$ with
$j(\alpha) \in X_c$, that is, $|\sigma_i(\alpha)| \leq c_i$ for all $i$. Each
complex embedding occurs in $\Nm(\alpha)$ together with its conjugate, and
$|\overline{\sigma}_i(\alpha)| = |\sigma_i(\alpha)|$, so
\[
  |\Nm(\alpha)|
  = \prod_{i=1}^{r_1} |\sigma_i(\alpha)|
    \prod_{i=r_1+1}^{r_1+r_2} |\sigma_i(\alpha)|^2
  \leq \prod_{i=1}^{r_1} c_i \prod_{i=r_1+1}^{r_1+r_2} c_i^2
  = M_K \Nm(I)
\]
by \eqref{eq:product}.
\end{proof}

\begin{proposition}\label{prop:small-ideal}
Every class $c \in \Cl_K$ contains an ideal $J \subseteq \OK$ with
$\Nm(J) \leq M_K$.
\end{proposition}

\begin{proof}
Let $c \in \Cl_K$ and write $c^{-1} = [\mathfrak{b}]$ for a fractional ideal
$\mathfrak{b}$. By Definition~\ref{def:classgroup} there is a nonzero
$\delta \in \OK$ with $I = \delta\mathfrak{b} \subseteq \OK$, and
$[I] = [\mathfrak{b}] = c^{-1}$ because $\delta\OK$ is principal. By
Lemma~\ref{lem:minkowski-bound} there is a nonzero $\alpha \in I$ with
$|\Nm(\alpha)| \leq M_K \Nm(I)$. Since $\alpha\OK \subseteq I$,
Fact~\ref{fact:dedekind} gives $\alpha\OK = IJ$ for an ideal
$J \subseteq \OK$. By Fact~\ref{fact:norm},
\[
  \Nm(J) = \frac{\Nm(\alpha\OK)}{\Nm(I)} = \frac{|\Nm(\alpha)|}{\Nm(I)}
  \leq M_K.
\]
Finally, $\alpha\OK = IJ$ is principal, so $[I][J] = 1$ in $\Cl_K$ and
$[J] = [I]^{-1} = c$.
\end{proof}

\section{Proof of Theorem~\ref{thm:finite}}\label{sec:proof}

\begin{lemma}\label{lem:finitely-many}
For every real number $B > 0$ there are only finitely many ideals
$J \subseteq \OK$ with $\Nm(J) \leq B$.
\end{lemma}

\begin{proof}
Let $J$ be such an ideal and $m = \Nm(J)$, a positive integer with
$m \leq B$. By Lemma~\ref{lem:norm-in-ideal}, $m \in J$, so
$m\OK \subseteq J$. By the correspondence theorem, the ideals of $\OK$
containing $m\OK$ are in bijection with the ideals of $\OK / m\OK$. By
Fact~\ref{fact:covolume}, $\OK \cong j(\OK) \cong \Z^n$ as abelian groups, so
$|\OK / m\OK| = m^n$ is finite, and $\OK/m\OK$ has finitely many ideals.
Hence, for each of the finitely many integers $1 \leq m \leq B$, there are
finitely many ideals $J$ with $\Nm(J) = m$.
\end{proof}

\begin{proof}[Proof of Theorem~\ref{thm:finite}]
By Proposition~\ref{prop:small-ideal}, every class in $\Cl_K$ contains one of
the ideals $J \subseteq \OK$ with $\Nm(J) \leq M_K$, and by
Lemma~\ref{lem:finitely-many} there are finitely many such ideals. Since
distinct classes are disjoint, the map sending each class to one such ideal
in it is injective, so $\Cl_K$ is finite.
\end{proof}

\begin{thebibliography}{9}
\bibitem{Mar} D.~A. Marcus, \emph{Number Fields}, Universitext,
  Springer-Verlag, New York, 1977.
\bibitem{Neu} J.~Neukirch, \emph{Algebraic Number Theory}, Grundlehren der
  mathematischen Wissenschaften 322, Springer-Verlag, Berlin, 1999.
\end{thebibliography}

\end{document}
